\documentclass[12pt, xcolor=dvipsnames,svgnames,x11names]{article}
\usepackage{xcolor}
\usepackage{xspace}
\usepackage{url}
\usepackage[colorlinks=true,bookmarks=false,linkcolor=black,urlcolor=blue,citecolor=black]{hyperref}
\usepackage{fancyhdr}
\usepackage[top=1in,bottom=1.2in,left=1in,right=1in]{geometry}
\usepackage{multicol}
\usepackage{pgfplots}
\usepackage{tikz}
\usetikzlibrary{circuits.logic.US, patterns}
\usepackage{mathpazo}
\usepackage{biolinum} 
\usepackage[all]{tcolorbox}
\usepackage{derivative}
\usepackage{../Lectures/rahul_math}
\renewcommand{\familydefault}{\sfdefault}

% --- Custom Color Palette from Image ---
\definecolor{headerRed}{HTML}{A3403D} % Dark red from image top
\definecolor{patternBg}{HTML}{F2D9D7} % Light pinkish bg from image bottom
\definecolor{binaryText}{HTML}{D1A7A4} % Muted color for binary digits

\usepackage{tikz}
\usetikzlibrary{positioning, arrows.meta, shapes.geometric, calc}

% --- Custom Colors ---
\definecolor{lstmBackground}{RGB}{240, 240, 255}
\definecolor{gateGreen}{RGB}{200, 255, 200}
\definecolor{matrixBlue}{RGB}{135, 206, 250}

% --- Background Pattern Setup ---
\usepackage[pages=all]{background}
\backgroundsetup{
scale=1,
color=black,
opacity=1,
angle=0,
contents={%
  \begin{tikzpicture}[remember picture,overlay]
    % Top Header Bar
    \fill[headerRed] (current page.north west) rectangle ([yshift=-1.2cm]current page.north east);
    % Bottom Binary Background
    \fill[patternBg] (current page.south west) rectangle ([yshift=3.5cm]current page.south east);
    % Decorative Line (from image)
    \draw[headerRed, line width=2pt] ([xshift=1in, yshift=-2.5cm]current page.north west) -- ([xshift=8in, yshift=-2.5cm]current page.north west);
    % Binary Text Overlay (Sample)
    \node[anchor=south, binaryText, font=\ttfamily\scriptsize, inner sep=0pt, yshift=1.5cm] at (current page.south) {
        \begin{tabular}{c@{\hspace{0.5em}}c@{\hspace{0.5em}}c@{\hspace{0.5em}}c@{\hspace{0.5em}}c@{\hspace{0.5em}}c@{\hspace{0.5em}}c}
        0 0 0 0 1 1 1 1 0 0 0 0 1 1 0 1 1 1 0 0 0 1 \\
        1 0 0 1 0 1 0 0 0 1 0 0 1 1 0 0 1 0 1 1 \\
        0 0 1 1 0 0 1 0 1 1 1 0 1 1 0 1 0 1 0 0 0 \\
        The University of Alabama in Huntsville
        \end{tabular}
    };
  \end{tikzpicture}
}
}

\newcommand{\hw}{Classwork 07//KL Divergence in DDPM}
\newcommand{\duedate}{March 23, 2026}

\usetikzlibrary{positioning, arrows.meta, shadows}

% --- Header Styling ---
\makeatletter
\renewcommand{\maketitle}{%
    \vspace*{0.2cm}
    \begin{center}
        \begin{tcolorbox}[
            enhanced,
            colback=white,
            colframe=headerRed,
            arc=0mm,
            title={\textcolor{white}{\bfseries \@title}},
            coltitle=white,
            fonttitle=\bfseries\large,
            attach boxed title to top left={xshift=10pt, yshift=-\tcboxedtitleheight/2},
            boxed title style={sharp corners, colback=headerRed, frame hidden}
        ]
        \large
        \vspace{10pt}
\noindent
\begin{tabular}{@{}p{3.5cm}p{6cm}r@{}}
\textbf{Student Name:} & \underline{\hspace{6cm}} & \textbf{Points:} 30 \\
\end{tabular}

\vspace{10pt}

\noindent
\begin{tabular}{@{}p{3.5cm}p{6cm}r@{}}
\textbf{A Number:} & \underline{\hspace{6cm}} & \textbf{Due Date:} \duedate \\
\end{tabular}
        \end{tcolorbox}
    \end{center}
    \vspace{10pt}
}
\makeatother



\pagestyle{fancy}
\fancyhf{}
\lhead{\small CPE 487/587, Spring 2026}
\rhead{\small \textit{\hw}}
\cfoot{\thepage}
\renewcommand{\headrulewidth}{0pt}

\usepackage{booktabs} % For professional lines


% Define a very light gray for the rows
\definecolor{tablegray}{gray}{0.95}

\begin{document}


\title{\hw}
\maketitle


KL Divergence for Denoising Diffusion Model (DDPM) is given by 

\begin{align*}
\text{KL}(q, p) & = \underbrace{\Ebb_{q(\xbf_0,\xbf_1, \cdots,\xbf_T)}\bigg[ -\log p(\xbf_T) - \sum_{t=1}^T\log \cfrac{p_\theta(\xbf_{t-1}| \xbf_t)}{q(\xbf_{t}| \xbf_{t-1})} \bigg]}_{=L} + \text{constant}
\end{align*}

Show that $L$ can be written as 

\begin{align*}
    \Ebb_{q(\xbf_0,\xbf_1, \cdots,\xbf_T)}\bigg[KL( q(\xbf_T | \xbf_0) || p(\xbf_T) ) + \sum_{t > 1} KL( q(\xbf_{t-1} | \xbf_t, \xbf_0) || p_\theta(\xbf_{t-1} | \xbf_t)) - \log p_\theta(\xbf_0 | \xbf_1) \bigg]
    \end{align*}
    
  Use the fundamental definition of KL divergence for proving.
  
  
  \subsection*{Proof:}
  
  The logarithm identity can be used for proving. First, we note that $\log\frac{a}{b} = \log a - \log b$. Use $\xbf_0,\xbf_1, \cdots,\xbf_T = \xbf_0:\xbf_T$ notation.
  
 

We also note that $KL(q,p) = \Ebb_q[ \log \frac{q}{p} ]$.  
  
  Then, 
  
  \begin{align*}
  L & = \Ebb_{q( \xbf_0:\xbf_T)}\bigg[ -\log p(\xbf_T) - \sum_{t=1}^T\log \cfrac{p_\theta(\xbf_{t-1}| \xbf_t)}{q(\xbf_{t}| \xbf_{t-1})} \bigg]\\
  & = \Ebb_{q( \xbf_0:\xbf_T)}\bigg[ -\log p(\xbf_T) - \sum_{t=1}^T\log p_\theta(\xbf_{t-1}| \xbf_t) + \sum_{t=1}^T\log q(\xbf_{t}| \xbf_{t-1}) \bigg]
  \end{align*}
  
  
   Also, we know from Bayes' theorem that
  
  \begin{align*}
  q(\xbf_t|\xbf_{t-1})  = \frac{q(\xbf_{t-1} | \xbf_t, \xbf_0) q(\xbf_t |\xbf_0)}{q(\xbf_{t-1}|\xbf_0)}
  \end{align*}
  Thus,
  
  \begin{align*}
  \sum_{t=}^T \log   q(\xbf_t|\xbf_{t-1}) = \log   q(\xbf_1|\xbf_0) + \sum_{t=2}^T \bigg[ \log q(\xbf_{t-1} | \xbf_t, \xbf_0) + \log q(\xbf_t | \xbf_0) - \log q(\xbf_{t-1} | \xbf_0) \bigg]
  \end{align*}
  
  If we expand the sum, then consecutive $\log q(\xbf_t | \xbf_0)$ with $+$ and $-$ will cancel out.
  
  Hence, 
  
  
    \begin{align*}
  \sum_{t=}^T \log   q(\xbf_t|\xbf_{t-1}) =\sum_{t=2}^T \log q(\xbf_{t-1} | \xbf_t, \xbf_0) + \log q(\xbf_T | \xbf_0)
  \end{align*}
  
  Then
  
  \begin{align*}
   L = \Ebb_q(\xbf_0:\xbf_T)\bigg [-\log p(\xbf_T) +  \sum_{t=2}^T \log q(\xbf_{t-1} | \xbf_t, \xbf_0) + \log q(\xbf_T | \xbf_0) -  \sum_{t=1}^T\log p_\theta(\xbf_{t-1}| \xbf_t)\bigg]
  \end{align*}
  
  
  Split the sum over $p_\theta$, we have
  
  \begin{align*}
  -\sum_{t=1}^T \log p_\theta(\xbf_{t-1}| \xbf_t) = -\log p_\theta(\xbf_0| \xbf_1) - \sum_{t=2}^T \log p_\theta(\xbf_{t-1}  | \xbf_t)
  \end{align*}
  
  Thus
  
  \begin{align*}
  L & = \Ebb_q(\xbf_0:\xbf_T)\bigg [-\log p(\xbf_T) +  \sum_{t=2}^T \log q(\xbf_{t-1} | \xbf_t, \xbf_0) + \log q(\xbf_T | \xbf_0) \\
  & -\log p_\theta(\xbf_0| \xbf_1) - \sum_{t=2}^T \log p_\theta(\xbf_{t-1}  | \xbf_t)\bigg]\\
  L & =  \Ebb_q(\xbf_0:\xbf_T)\bigg [\frac{ \log q(\xbf_T | \xbf_0)}{\log p(\xbf_T)} + \sum_{t=2}^T \log \frac{q(\xbf_{t-1}|\xbf_t, \xbf_0}{p_\theta(\xbf_{t-1} | \xbf_t)} - \log p_\theta(\xbf_0 | \xbf_1) \bigg]
  \end{align*}
  
  Each $\log$ is argument to KL as per its definition, thus, taking expectations over $q(\xbf_0 : \xbf_T)$
  
  \begin{align*}
  L = \Ebb_q(\xbf_0: \xbf_T) \bigg[ KL(q(\xbf_T | \xbf_0)|| p(\xbf_T)) + \sum_{t=2}^T KL(q(\xbf_{t-1} | \xbf_t, \xbf_0) || p_\theta(\xbf_{t-1} | \xbf_t)) - \log p_\theta(\xbf_0 | \xbf_1) \bigg]
  \end{align*}
  

\end{document}